\chapter{Chapter05 — Linear Threshold Functions}
%%% generated by the blueprint pipeline from FABL.Chapter05.LinearThresholdFunctions
\section{Overview}

Declarations of the Lean module \texttt{FABL.Chapter05.LinearThresholdFunctions}.

\section{Declarations}

\begin{theorem}[Fourier coefficient bounded by the first norm]
\lean{FABL.abs_fourierCoeff_le_uniformLpNorm_one}
\label{FABL.abs_fourierCoeff_le_uniformLpNorm_one}
\leanok
\uses{FABL.fourierCoeff, FABL.monomial, FABL.monomial_sq, FABL.uniformLpNorm}
\difficulty{4}
Let $n$ be a nonnegative integer and let $f\colon\{\pm1\}^n\to\mathbb R$ be any
function. For every subset $S$ of the coordinates, the absolute value of the Fourier
coefficient of $f$ on $S$ is at most its uniform $L^1$ norm.
\end{theorem}

\begin{theorem}[Fourier one norm bounds pointwise values]
\lean{FABL.abs_apply_le_fourierOneNorm}
\label{FABL.abs_apply_le_fourierOneNorm}
\leanok
\uses{FABL.fourierCoeff, FABL.fourierOneNorm, FABL.fourier_expansion, FABL.monomial, FABL.monomial_sq}
\difficulty{4}
Let $n$ be a natural number, and let $f$ be a real-valued function on the Boolean cube
$\{-1,1\}^n$. At every point $x$ of the cube, the absolute value $|f(x)|$ is at most the
sum of the absolute values of all Fourier coefficients of $f$.
\end{theorem}

\begin{definition}[Polynomial threshold representation]
\lean{FABL.IsPolynomialThresholdRepresentation}
\statementsource{Analysis-of-Boolean-Functions-by-Ryan-ODonnell}{pdf-ffbe9c7b3306-p115-b005}
\label{FABL.IsPolynomialThresholdRepresentation}
\leanok
\uses{FABL.BooleanFunction, FABL.thresholdSign}
Let $n$ be a natural number, let $f$ map the $n$-dimensional Boolean cube $\\{-1,1\\}^n$
to the signs, and let $p:\\{-1,1\\}^n\\to\\mathbb R$. The function $f$ is represented by
the polynomial threshold of $p$ if, at every input $x$, its value is the threshold sign
of $p(x)$.
\end{definition}

\begin{definition}[Polynomial threshold function]
\lean{FABL.IsPolynomialThreshold}
\statementsource{Analysis-of-Boolean-Functions-by-Ryan-ODonnell}{pdf-ffbe9c7b3306-p115-b005}
\label{FABL.IsPolynomialThreshold}
\leanok
\uses{FABL.BooleanFunction, FABL.IsPolynomialThresholdRepresentation, FABL.fourierDegree}
Let $n$ and $k$ be natural numbers, and let $f$ map the $n$-dimensional Boolean cube to
the signs. Then $f$ is a polynomial threshold function of degree at most $k$ if there is
a real-valued polynomial on the cube whose threshold sign agrees with $f$ at every input
and whose Fourier degree is at most $k$.
\end{definition}

\begin{definition}[Polynomial Fourier sparsity]
\lean{FABL.polynomialSparsity}
\statementsource{Analysis-of-Boolean-Functions-by-Ryan-ODonnell}{pdf-ffbe9c7b3306-p116-b002}
\label{FABL.polynomialSparsity}
\leanok
\uses{FABL.fourierSupport}
For a real-valued function $p$ on the $n$-dimensional Boolean cube, define its
polynomial sparsity to be the cardinality of its Fourier support.
\end{definition}

\begin{definition}[Four-bit equality function]
\lean{FABL.fourBitEquality}
\statementsource{Analysis-of-Boolean-Functions-by-Ryan-ODonnell}{
  pdf-ffbe9c7b3306-p034-b003,
  pdf-ffbe9c7b3306-p115-b006
}
\label{FABL.fourBitEquality}
\leanok
\uses{FABL.BooleanFunction}
Define a function on four signs by assigning $1$ when all four coordinates are equal,
and $-1$ otherwise.
\end{definition}

\begin{definition}[Four-bit equality polynomial]
\lean{FABL.fourBitEqualityPolynomial}
\statementsource{Analysis-of-Boolean-Functions-by-Ryan-ODonnell}{pdf-ffbe9c7b3306-p115-b006}
\label{FABL.fourBitEqualityPolynomial}
\leanok
\uses{FABL.monomial}
For a sign vector $x=(x_0,x_1,x_2,x_3)\\in\{-1,1\}^4$, define $p(x)=-3+\sum_{0\le
i<j<4}x_i x_j$.
\end{definition}

\begin{definition}[Quadratic frequencies for four-bit equality]
\lean{FABL.fourBitEqualityQuadraticFrequencies}
\statementsource{Analysis-of-Boolean-Functions-by-Ryan-ODonnell}{pdf-ffbe9c7b3306-p115-b006}
\label{FABL.fourBitEqualityQuadraticFrequencies}
\leanok
Define the set of quadratic frequencies to be the six two-element subsets of
$\{0,1,2,3\}$.
\end{definition}

\begin{definition}[Coefficients of the four-bit polynomial]
\lean{FABL.fourBitEqualityPolynomialCoefficients}
\statementsource{Analysis-of-Boolean-Functions-by-Ryan-ODonnell}{pdf-ffbe9c7b3306-p115-b006}
\label{FABL.fourBitEqualityPolynomialCoefficients}
\leanok
\uses{FABL.fourBitEqualityQuadraticFrequencies}
For each subset $S$ of $\{0,1,2,3\}$, define its coefficient to be $-3$ if $S$ is empty,
$1$ if $S$ is one of the six two-element subsets, and $0$ otherwise.
\end{definition}

\begin{theorem}[Expansion of the four-bit equality polynomial]
\lean{FABL.fourBitEqualityPolynomial_expansion}
\label{FABL.fourBitEqualityPolynomial_expansion}
\leanok
\uses{FABL.fourBitEqualityPolynomial, FABL.fourBitEqualityPolynomialCoefficients, FABL.fourBitEqualityQuadraticFrequencies, FABL.monomial, FABL.multilinearPolynomial}
\difficulty{4}
For every sign vector $x\in\{−1,1\}^4$, the four-bit equality polynomial equals the
multilinear polynomial whose coefficient is $-3$ for the empty set, $1$ for each
two-element subset of the four coordinates, and $0$ for every other subset.
\end{theorem}

\begin{theorem}[Four-bit equality polynomial coefficients]
\lean{FABL.fourierCoeff_fourBitEqualityPolynomial}
\label{FABL.fourierCoeff_fourBitEqualityPolynomial}
\leanok
\uses{FABL.fourBitEqualityPolynomial, FABL.fourBitEqualityPolynomialCoefficients, FABL.fourBitEqualityPolynomial_expansion, FABL.fourierCoeff, FABL.fourier_expansion_unique}
\difficulty{1}
On the sign cube $\{\pm1\}^4$, define $P(x)=-3+\sum_{0\le j<k<4}x_jx_k$. For every
subset $S\subseteq\{0,1,2,3\}$, the Fourier coefficient of $P$ at $S$ is $-3$ if $S$ is
empty, $1$ if $S$ has two elements, and $0$ otherwise.
\end{theorem}

\begin{theorem}[Polynomial threshold representation of four-bit equality]
\lean{FABL.fourBitEquality_polynomialThresholdRepresentation}
\proofsource{Analysis-of-Boolean-Functions-by-Ryan-ODonnell}{pdf-ffbe9c7b3306-p115-b006}
\label{FABL.fourBitEquality_polynomialThresholdRepresentation}
\leanok
\uses{FABL.IsPolynomialThresholdRepresentation, FABL.fourBitEquality, FABL.fourBitEqualityPolynomial, FABL.monomial, FABL.thresholdSign}
\difficulty{4}
For an input $x\in\{-1,1\}^4$, let $f(x)=-1$ when all four coordinates are equal and
$f(x)=1$ otherwise. Let $p(x)=-3+\sum_{0\le i<j<4}x_i x_j$. Then $f(x)=-1$ when $p(x)<0$
and $f(x)=1$ when $p(x)\ge 0$ for every $x\in\{-1,1\}^4$.
\end{theorem}

\begin{theorem}[Degree of the four-bit equality polynomial]
\lean{FABL.fourierDegree_fourBitEqualityPolynomial_le}
\label{FABL.fourierDegree_fourBitEqualityPolynomial_le}
\leanok
\uses{FABL.fourBitEqualityPolynomial, FABL.fourBitEqualityPolynomialCoefficients, FABL.fourBitEqualityQuadraticFrequencies, FABL.fourierCoeff_fourBitEqualityPolynomial, FABL.fourierDegree, FABL.fourierDegree_le_iff}
\difficulty{5}
On $\{-1,1\}^4$, define $P(x)=-3+\sum_{0\le i<j<4}x_i x_j$. The Fourier degree of $P$ is
at most $2$.
\end{theorem}

\begin{theorem}[Degree-two threshold representation of four-bit equality]
\lean{FABL.fourBitEquality_isPolynomialThreshold}
\proofsource{Analysis-of-Boolean-Functions-by-Ryan-ODonnell}{pdf-ffbe9c7b3306-p115-b006}
\label{FABL.fourBitEquality_isPolynomialThreshold}
\leanok
\uses{FABL.IsPolynomialThreshold, FABL.fourBitEquality, FABL.fourBitEqualityPolynomial, FABL.fourBitEquality_polynomialThresholdRepresentation, FABL.fourierDegree_fourBitEqualityPolynomial_le}
\difficulty{1}
The Boolean function on sign vectors of length four that is $1$ when all four
coordinates are equal and $-1$ otherwise has a polynomial threshold representation of
Fourier degree at most $2$.
\end{theorem}

\begin{theorem}[Fourier support of the four-bit equality polynomial]
\lean{FABL.fourierSupport_fourBitEqualityPolynomial}
\label{FABL.fourierSupport_fourBitEqualityPolynomial}
\leanok
\uses{FABL.fourBitEqualityPolynomial, FABL.fourBitEqualityPolynomialCoefficients, FABL.fourBitEqualityQuadraticFrequencies, FABL.fourierCoeff_fourBitEqualityPolynomial, FABL.fourierSupport, FABL.mem_fourierSupport}
\difficulty{4}
On the sign cube of dimension four, define $P(x)=-3+\sum_{0\le i<j<4}x_i x_j$. Its
Fourier support, meaning the subsets with nonzero Fourier coefficient, consists exactly
of the empty set and the six two-element subsets of $\{0,1,2,3\}$.
\end{theorem}

\begin{theorem}[Polynomial sparsity four bit equality polynomial]
\lean{FABL.polynomialSparsity_fourBitEqualityPolynomial}
\statementsource{Analysis-of-Boolean-Functions-by-Ryan-ODonnell}{
  pdf-ffbe9c7b3306-p115-b006,
  pdf-ffbe9c7b3306-p116-b002
}
\label{FABL.polynomialSparsity_fourBitEqualityPolynomial}
\leanok
\uses{FABL.fourBitEqualityPolynomial, FABL.fourierSupport_fourBitEqualityPolynomial, FABL.polynomialSparsity}
\difficulty{3}
% dataset-pipeline: stub -- prose pending build_informal_statements.py
Statement pending regeneration.
\end{theorem}

\begin{theorem}[Polynomial threshold approximation from noise sensitivity]
\lean{FABL.exists_polynomialThreshold_relativeHammingDist_le_three_mul_noiseSensitivity}
\proofsource{Analysis-of-Boolean-Functions-by-Ryan-ODonnell}{pdf-ffbe9c7b3306-p115-b008}
\label{FABL.exists_polynomialThreshold_relativeHammingDist_le_three_mul_noiseSensitivity}
\leanok
\uses{FABL.BooleanFunction, FABL.BooleanFunction.toReal, FABL.IsPolynomialThreshold, FABL.fourierCoeff_lowDegreePart, FABL.fourierDegree, FABL.fourierDegree_le_iff, FABL.fourierWeightAbove, FABL.fourierWeightAboveReal, FABL.isFourierSpectrumConcentratedUpTo_noiseSensitivity, FABL.lowDegreePart, FABL.noiseSensitivity, FABL.relativeHammingDist, FABL.relativeHammingDist_thresholdSign_le_of_uniformLpNorm_two_sq_le, FABL.thresholdSign, FABL.two_div_one_sub_exp_neg_two_mul_noiseSensitivity_le_three, FABL.uniformLpNorm, FABL.uniformLpNorm_sub_lowDegreePart_sq}
\difficulty{5}
Let $f:\\{−1,1\\}^n\\to\\{−1,1\\}$ and let $0<\\delta\\le1/2$. There is a Boolean
function $g$ with a polynomial threshold representation of Fourier degree at most
$\\lfloor1/\\delta\\rfloor$, such that $\\lfloor1/\\delta\\rfloor\\le1/\\delta$ and the
relative Hamming distance between $f$ and $g$ is at most three times the noise
sensitivity of $f$ at noise rate $\\delta$.
\end{theorem}

\begin{theorem}[Fourier coefficient mass of a threshold function]
\lean{FABL.one_le_sum_abs_fourierCoeff_of_polynomialThresholdRepresentation}
\proofsource{Analysis-of-Boolean-Functions-by-Ryan-ODonnell}{
  pdf-ffbe9c7b3306-p116-b006,
  pdf-ffbe9c7b3306-p116-b007
}
\label{FABL.one_le_sum_abs_fourierCoeff_of_polynomialThresholdRepresentation}
\leanok
\uses{FABL.BooleanFunction, FABL.BooleanFunction.toReal, FABL.IsPolynomialThresholdRepresentation, FABL.abs_fourierCoeff_le_uniformLpNorm_one, FABL.fourierCoeff, FABL.fourierSupport, FABL.mem_fourierSupport, FABL.plancherel, FABL.thresholdSign_of_neg, FABL.thresholdSign_of_nonneg, FABL.uniformInner, FABL.uniformLpNorm}
\difficulty{6}
Let $f:\\{−1,1\\}^n\to\{−1,1\}$ be a Boolean function, let $p:\\{−1,1\\}^n\to\mathbb R$
be a nonzero function, and let $\mathcal F$ be a family of subsets of $[n]$. Suppose
$f(x)$ is the sign of $p(x)$, with zero assigned the positive sign, for every input $x$,
and every set with a nonzero Fourier coefficient of $p$ belongs to $\mathcal F$. Then
the sum over $S\in\mathcal F$ of the absolute values of the Fourier coefficients of the
real-valued function associated to $f$ is at least $1$.
\end{theorem}
